% Standard edition. Literature checked through 25 September 2026.
% Normal shared typography; no manual page packing or forced page count.
\documentclass[12pt,reqno]{amsart}
\usepackage{math-review}
\newcommand{\ii}{\sqrt{-1}}
\newcommand{\ddc}{\ii\partial\bar\partial}
\newcommand{\R}{\mathbb R}
\newcommand{\C}{\mathbb C}
\newcommand{\Th}{\Theta_\omega}
\newcommand{\hth}{\widehat\theta}
\newcommand{\arccot}{\operatorname{arccot}}
\newcommand{\osc}{\operatorname{osc}}
\newcommand{\Ree}{\operatorname{Re}}
\newcommand{\Imm}{\operatorname{Im}}
\title[Scalar dHYM: solvability and mechanisms]{The Scalar Deformed Hermitian--Yang--Mills Equation:\\ Solvability and Its Mechanisms}
\author{}
\hypersetup{pdftitle={The Scalar Deformed Hermitian--Yang--Mills Equation: Solvability and Its Mechanisms},pdfsubject={English standard mathematical survey; literature cutoff 2026-09-25},pdfkeywords={dHYM, subsolution, numerical criterion, tangent phase flow}}
\date{Standard edition\quad---\quad25 September 2026}
\begin{document}
\begin{abstract}
The scalar deformed Hermitian--Yang--Mills equation asks for a representative of a real $(1,1)$ class with constant lifted phase. We explain smooth solvability in the positive supercritical range through three interfaces: a subsolution prevents loss of elliptic estimates; numerical positivity produces a subsolution; and, at hypercritical phase, the tangent phase flow constructs the solution. The surface Monge--Amp\`ere reduction supplies an explicit model, while higher-dimensional ray tests expose the importance of the real branch. Critical phases and non-strict positivity are treated as distinct boundaries, not as consequences of the strict theory.
\end{abstract}
\maketitle

\section{Introduction}\label{sec:intro}
Let $(X,\omega)$ be a compact connected K\"ahler manifold of complex dimension $n\geq2$, let $\alpha$ be a smooth closed real $(1,1)$ form, and write $A=[\alpha]$ and $\Omega=[\omega]$. The unknown is a real smooth potential $u$, modulo constants. Put
\begin{equation}\label{eq:equation}
 \alpha_u=\alpha+\ddc u,\qquad
 \Th(\alpha_u)=\sum_{j=1}^n\arctan\lambda_j(\omega^{-1}\alpha_u)=\hth.
\end{equation}
Each arctangent takes values in $(-\pi/2,\pi/2)$. Unless explicitly stated otherwise, the target is \emph{positive supercritical}:
\begin{equation}\label{eq:range}
 (n-2)\pi/2<\hth<n\pi/2,\qquad
 Z=\int_X(\omega+\ii\alpha)^n,\qquad e^{-\ii\hth}Z\in\R_{>0}.
\end{equation}
The last condition is necessary, but does not force representatives onto the specified real branch. Smooth solutions are unique modulo constants.

The principal analytic answer is an equivalence with a strict, pointwise condition on one representative. A \emph{$C$-subsolution} is a smooth $\underline u$ whose eigenvalues $\mu_j$ satisfy
\begin{equation}\label{eq:subsolution}
 \sum_{j\ne k}\arctan\mu_j>\hth-\pi/2
 \quad\text{at every }x\in X\text{ and for every }k=1,\ldots,n.
\end{equation}
\begin{theorem}\label{thm:analytic}
Under \eqref{eq:range}, equation \eqref{eq:equation} has a smooth solution if and only if a smooth $C$-subsolution exists. No additional lower bound on the full phase $\Th(\alpha_{\underline u})$ is required.
\end{theorem}
For $n\geq3$ the version used here is Sun \cite[Proposition 4.1]{Sun}, whose assumption even allows a $C^2$ subsolution; for $n=2$ it follows from the surface reduction of Jacob--Yau \cite[Theorem 1.2]{JY}. The foundational estimates are due to Collins--Jacob--Yau \cite{CJY}; their original existence theorem has an additional phase hypothesis. Section~\ref{sec:analytic} explains both the estimates and the later continuity argument, rather than silently attributing the stronger theorem to the earlier proof.

For a numerical answer, it is convenient to use the complementary phase
\begin{equation}\label{eq:Q}
 \theta=n\pi/2-\hth\in(0,\pi),\qquad
 Q_\omega(\beta)=\sum_j\arccot\lambda_j(\omega^{-1}\beta),
\end{equation}
where $\arccot\colon\R\to(0,\pi)$. Define the real cohomology class
\begin{equation}\label{eq:F}
 F_p(B,\Omega;\theta)=\Ree(B+\ii\Omega)^p
                  -\cot\theta\,\Imm(B+\ii\Omega)^p.
\end{equation}
Integrating over a singular analytic subvariety means pairing with its fundamental cycle.
\begin{theorem}\label{thm:numerical}
Under \eqref{eq:range}, smooth solvability is equivalent to the following tests for every real $t\geq0$ and every irreducible analytic subvariety $V\subseteq X$ of dimension $1\leq p\leq n$:
\begin{equation}
 \int_V F_p(A+t\Omega,\Omega;\theta)
 \begin{cases}>0,&p<n,\\ \geq0,&p=n.\end{cases}
\end{equation}
The top integral is zero at $t=0$ by phase compatibility. If $X$ is projective, the strict inequalities for all proper positive-dimensional $V$ at $t=0$ alone suffice; the same top-degree compatibility is still required.
\end{theorem}
This is the ray specialization and projective consequence of Chu--Lee--Takahashi \cite[Theorem 1.3 and Corollaries 1.4--1.5]{CLT}. The projective assertion does not require $A$ or $\Omega$ to be rational. The full ray, including its top-dimensional test, cannot simply be discarded on a general compact K\"ahler manifold.

Two concrete consequences organize the explanations. On a surface with $s=A\cdot\Omega>0$, set
\begin{equation}\label{eq:surface-data}
 c=\frac{\Omega^2-A^2}{2s},\qquad B=A+c\Omega.
\end{equation}
Then smooth solvability on the phase in $(0,\pi)$ is equivalent to $B$ being a K\"ahler class, or equivalently to $B\cdot C>0$ for every irreducible curve $C$; the positive-cone conditions are automatic for this $B$. In higher dimension, when $\hth>(n-1)\pi/2$, a smooth subsolution also ensures that the tangent phase flow from \emph{every} smooth initial potential satisfying $|\Th(\alpha_{u_0})-\hth|<\pi/2$ exists for all time and converges smoothly to a solution \cite[Theorem 1.4]{Tak}. The initial potential need not itself have hypercritical phase.

The scope is scalar smooth solvability on a fixed compact K\"ahler background. The surface calculation, the analytic closure, the numerical construction and this one evolution receive detailed attention. Properness criteria, a parallel development of other flows, and categorical stability are not needed for these arguments. References are checked at the displayed versions through the date above; the recent critical-phase theorem in Section~\ref{sec:boundary} is explicitly a preprint result. The source check concerns statements and selected proof interfaces, not a complete independent verification of all cited proofs.

\section{Phase and the surface model}\label{sec:surface}
\subsection{What the scalar equation fixes}
Diagonalization gives the pointwise identity
\begin{equation}\label{eq:det}
 \frac{(\omega+\ii\alpha_u)^n}{\omega^n}
   =\prod_j(1+\ii\lambda_j)
   =\prod_j\sqrt{1+\lambda_j^2}\,e^{\ii\Th(\alpha_u)}.
\end{equation}
Consequently \eqref{eq:equation} implies \eqref{eq:range}. The equation that the rotated imaginary part vanishes fixes only a phase modulo $\pi$; positivity of its real part improves this to modulo $2\pi$, not to a specified real lift. The sum in \eqref{eq:equation}, rather than an arbitrary logarithm of a determinant, fixes the branch. Similarly $Q_\omega=n\pi/2-\Th$ is a real-valued identity, not a congruence.

The linearization is elliptic at every finite Hermitian matrix:
\begin{equation}\label{eq:linearization}
 D\Th\big|_{\alpha_u}(v)
      =\sum_j\frac{v_{j\bar j}}{1+\lambda_j^2}
\end{equation}
in an $\omega$-unitary diagonal frame. For two smooth solutions with the same target, integrate this linearization along their affine segment. The resulting uniformly elliptic linear equation for their difference and the maximum principle prove uniqueness modulo constants. This argument does not assert concavity of $\Th$ on that segment.

No Calabi--Yau metric is required. For a line bundle one may fix the convention $\alpha=\ii F_h/(2\pi)$, $F_h=-\partial\bar\partial\log h$; replacing $h$ by $e^{-2\pi u}h$ gives $\alpha_u$. The scalar problem also makes sense for nonintegral real classes. Locally, the analogous determinant of $I+\ii D^2f$ is the phase of a Lagrangian graph in $\C^n$. This explains the special Lagrangian terminology, but is not a global identification of solutions or of flows; see \cite{JY}. No Thomas--Yau or categorical equivalence is used below.

\subsection{A complete two-way reduction}
On a surface the imaginary part of $Z$ is $2s>0$, so
\[
 \hth=\operatorname{Arg}(\Omega^2-A^2+2\ii s)\in(0,\pi),
 \qquad c=\cot\hth.
\]
Expanding the rotated equation gives
\begin{equation}\label{eq:surface-expand}
 \alpha_u^2+2c\,\alpha_u\wedge\omega-\omega^2=0.
\end{equation}
For a genuine solution, each eigenangle exceeds $\hth-\pi/2$. Hence $\lambda_j+c>0$ and the shifted form $\beta=\alpha_u+c\omega$ satisfies
\begin{equation}\label{eq:MA}
 \beta>0,\qquad \beta^2=(1+c^2)\omega^2,\qquad [\beta]=B.
\end{equation}
This proves the necessity of $B$ being K\"ahler, not merely nef.

Conversely, direct intersection algebra gives
\begin{equation}\label{eq:volume}
 B^2=A^2+2cs+c^2\Omega^2=(1+c^2)\Omega^2.
\end{equation}
If $B$ is K\"ahler, Yau's complex Monge--Amp\`ere theorem \cite{Yau} therefore supplies a smooth positive representative satisfying \eqref{eq:MA}. The $\partial\bar\partial$ lemma writes it as $\alpha+c\omega+\ddc u$. This gives \eqref{eq:surface-expand}, but one must still select the phase. Writing $\mu_j=\lambda_j+c>0$, the determinant equation gives $\mu_1\mu_2=1+c^2$, and therefore
\begin{equation}\label{eq:branch-back}
 \lambda_1+\lambda_2\geq2\sqrt{1+c^2}-2c>0.
\end{equation}
Thus $(\omega+\ii\alpha_u)^2/\omega^2$ lies in the upper half-plane. Its phase is the prescribed $\hth\in(0,\pi)$, rather than the other solution of the equation modulo $\pi$. This completes the converse, including the branch issue left invisible by squaring. It also extends the line-bundle formulation of \cite{JY} to any real $(1,1)$ class without an integrality assumption.

\subsection{Numerical strictness and its boundary}
To apply the surface K\"ahler-cone criterion \cite{DP}, it remains to place $B$ in the correct positive component. Put $a=A^2$ and $v=\Omega^2$. The Hodge index inequality $av\leq s^2$ gives
\[
 B\cdot\Omega=s+cv
   =\frac{2s^2+v^2-av}{2s}>0,
 \qquad B^2=(1+c^2)v>0.
\]
In this component, strict positive intersection with every irreducible curve is equivalent to being K\"ahler. Thus the numerical assertion in the introduction includes, rather than suppresses, the positive-cone qualification. For $n=2$, condition \eqref{eq:subsolution} is exactly positivity of $\alpha_{\underline u}+c\omega$, proving Theorem~\ref{thm:analytic} in this dimension.

A concrete boundary prevents a misleading replacement of $>$ by $\geq$. On the blow-up of $\mathbb P^2$ at a point, let $H$ be the pullback hyperplane class and $E$ the exceptional curve. Take
\[
 \Omega=\frac{2H-E}{\sqrt3},\qquad A=H.
\]
Here $H^2=1$, $E^2=-1$, $H\cdot E=0$, and $\Omega$ is K\"ahler. We have $A^2=\Omega^2=1$, $s=2/\sqrt3$, $c=0$ and $\hth=\pi/2$. The class $B=H$ is nef and big, but $B\cdot E=0$. A smooth dHYM solution would give a K\"ahler form in $H$ whose integral over $E$ is both positive and zero, a contradiction. This example concerns failure of \emph{smooth} solvability; it does not rule out a separately defined weak solution.

\section{Why a subsolution closes the elliptic problem}\label{sec:analytic}
\subsection{Preventing escape along a phase level}
Fix a point and the eigenvalues $\mu$ of a subsolution. If $\lambda\geq\mu$ coordinatewise and $\sum\arctan\lambda_j=\hth$, no coordinate can tend to $+\infty$: otherwise the limiting phase is at least
\[
 \pi/2+\sum_{j\ne k}\arctan\mu_j>\hth.
\]
The strict gap is uniform on compact $X$. This is the geometric content of a bounded upper-orthant slice of a phase level. It is stronger than checking a sign after integration, yet much weaker than finding a solution. A solution itself satisfies the strict inequalities because every omitted angle is strictly less than $\pi/2$.

The quantitative form uses convexity of supercritical phase superlevel sets. Outside a fixed bounded region, the normal to the solution level either has a definite component toward the subsolution or all its coordinate components are comparable. In the linearized maximum-principle calculation, this yields either a positive contribution from the subsolution barrier or sufficient ellipticity in every direction. Collins--Jacob--Yau \cite[Proposition 3.5]{CJY} establish the version needed for their operator. The relevant constants depend on the strict subsolution gap; a non-strict inequality does not supply them.

\subsection{The estimates and the missing regularity interface}
The estimate package in \cite[Sections 3--6]{CJY} has a definite order. Normalize $\sup_Xu=0$ and absorb the subsolution into the reference. A fixed supercritical gap implies a lower eigenvalue bound and hence $\Delta_\omega u\geq-C$; the Green function gives $\int_X|u|\omega^n\leq C$. An integral bound does not yet control the depth of a narrow minimum.

The Alexandrov--Bakelman--Pucci argument prevents precisely that phenomenon. In a coordinate ball centered at a minimum, perturb $u$ by a small quadratic function. On its lower contact set the real Hessian of the perturbed function is nonnegative. Its complex Hessian therefore places $\alpha_u$ above a small downward perturbation of the subsolution. Strictness keeps the corresponding upper-orthant phase slices uniformly bounded, so the eigenvalues of $\alpha_u$ are bounded on that contact set. Positivity of the real Hessian and its bounded trace then bound its real determinant. Schematically, for a fixed small quadratic parameter $\varepsilon$ and the contact set $\mathcal C$,
\[
 c\varepsilon^{2n}\leq
  \int_{\mathcal C}\det_{\R}D^2(u+\varepsilon|z|^2)
  \leq C\,\operatorname{Vol}(\mathcal C).
\]
Thus $\mathcal C$ has a uniform positive volume. By the contact-set construction, $u\leq\inf_Xu+C\varepsilon$ there. The global $L^1$ bound now forces $-\inf_Xu\leq C$. The fixed coordinate comparisons and contact-set inequalities are those of \cite[Proposition 3.6]{CJY}; the point is that bounded level slices enter before any global Hessian estimate. This proves a bound for the potential without requiring the reference form to be K\"ahler.

The maximum principle for a largest-eigenvalue test function, with gradient and subsolution corrections, next gives
\begin{equation}\label{eq:hessian-grad}
 \|\ddc u\|_{C^0}\leq C\bigl(1+\|\nabla u\|_{C^0}^2\bigr).
\end{equation}
The normal alternative above supplies the favorable term when a Hessian eigenvalue becomes large. The constants use fixed background data, a positive supercritical gap and the subsolution. Formula \eqref{eq:hessian-grad} is not yet a closed second-order bound.

Here is why the gradient cannot diverge. If $M_j=|\nabla u_j(x_j)|\to\infty$, rescale space by $M_j^{-1}$ around $x_j$, leaving the amplitude of $u_j$ unchanged. A supercritical phase with fixed gap implies a uniform lower bound on every eigenvalue: subtracting the other $n-1$ angles bounds the remaining angle away from $-\pi/2$. In rescaled coordinates this gives
\[
 \ddc v_j\geq-O(M_j^{-2})\omega_{\rm Eucl},\qquad
 |\ddc v_j|\leq C,
 \qquad |\nabla v_j(0)|=1.
\]
The $C^0$ bound persists on expanding balls. Local elliptic compactness, using the Laplacian bound, produces a bounded plurisubharmonic limit on $\C^n$ with nonzero gradient at the origin. Restriction to complex lines and the bounded subharmonic Liouville theorem force it to be constant. This contradiction closes the gradient estimate and then \eqref{eq:hessian-grad} \cite[Proposition 5.1]{CJY}.

Bounded eigenvalues now imply uniform ellipticity, but \emph{not by themselves} a $C^{2,\gamma}$ estimate. Indeed
\[
 \frac{d^2}{d\lambda^2}\arctan\lambda
       =-\frac{2\lambda}{(1+\lambda^2)^2}
\]
is positive at negative $\lambda$, which can occur at supercritical phase. Invoking concavity of the original phase operator here would be incorrect. Section 6 of \cite{CJY} instead uses the convex geometry of its level sets to obtain a suitable concave elliptic defining equation. On the controlled Hessian range one passes through the Hermitian projection to a real uniformly elliptic concave framework. The associated regularity and Liouville argument supplies $C^{2,\gamma}$ control; Schauder estimates then give all higher orders. This replacement of the operator, not a change of the equation's solution set, is the regularity interface needed for closedness.

\subsection{Which continuity theorem is being proved?}
The original existence theorem \cite[Theorem 1.2]{CJY} assumes, in addition to \eqref{eq:subsolution}, either
\[
 \Th(\alpha_{\underline u})>(n-2)\pi/2
 \quad\text{or}\quad
 \hth\geq(n-2+2/n)\pi/2.
\]
The second alternative implies the first by summing the $n$ strict subsolution inequalities. It is not legitimate to delete both alternatives while retaining only that paper's continuity argument. Later work removes the extra restriction; the implementation below follows Sun's alternative proof \cite[Proposition 4.1]{Sun}. 

Absorb the subsolution into the reference form, so $\underline u=0$, and use the notation \eqref{eq:Q}--\eqref{eq:F}. Set
\begin{equation}\label{eq:twisted}
 F_n(\alpha+t\omega+\ddc u_t,\omega;\theta)=c_t\omega^n,
 \qquad
 c_t=\frac{\int_X F_n(A+t\Omega,\Omega;\theta)}{\int_X\omega^n}.
\end{equation}
Here $c_0=0$. The lower-degree positivity implied by the subsolution makes $c_t>0$ for $t>0$: differentiation of the numerator gives $n\int_X F_{n-1}(\alpha+t\omega,\omega;\theta)\wedge\omega>0$. For large $t$, the reference is strongly positive and Chen's twisted-equation result, in the form used by Sun, starts the path. It does \emph{not} mean that $u_t=0$ solves a generally nonconstant pointwise right-hand side.

The path is considered in the controlled cone
\[
 P(\lambda):=\max_k\sum_{j\ne k}\arccot\lambda_j<\theta,
 \qquad Q(\lambda)<\Theta_0<\pi,
\]
with $\theta<\Theta_0$. Writing
\[
 D(\lambda)=\Imm\prod_j(\lambda_j+\ii)
            =\sin Q(\lambda)\prod_j\sqrt{1+\lambda_j^2}>0,
\]
its operator can be expressed as
\begin{equation}\label{eq:twisted-operator}
 G_t(\lambda)=\cot Q(\lambda)-\frac{c_t}{D(\lambda)},
 \qquad G_t(\lambda)=\cot\theta.
\end{equation}
The ellipticity and concavity available for this twisted equation are imported from Chen's theory, with its cone and right-hand-side restrictions, not asserted for arbitrary Hermitian matrices \cite[Proposition 2.8]{CLT}.

The crucial uniform barrier uses the \emph{partial} phase, not the full phase of $\alpha$. Strictness permits a small $\varepsilon>0$ such that $P(\alpha-\varepsilon\omega)$ is still below $\theta$ by a fixed margin. In a unitary frame increase one eigenvalue of $\alpha-\varepsilon\omega$ by $N$; compression interlacing gives the same limiting bound for a boost in any unitary direction. As $N\to\infty$, its arccotangent tends to zero, so its full phase approaches an $(n-1)$-angle sum strictly below $\theta$. Its determinant density $D$ grows linearly in $N$, with positive leading coefficient. Since $c_t$ is bounded on a fixed interval $[0,T]$, the second term in \eqref{eq:twisted-operator} tends uniformly to zero. For one sufficiently large $N$, the boosted form consequently satisfies
\[
 G_t(\hbox{boosted form})>\cot\theta+\xi
\]
for some fixed $\xi>0$, in every boosted direction and for all $t\in[0,T]$.

The conversion to an estimate is quantitative. At a diagonal solution matrix let $g_j=\partial G_t/\partial\lambda_j>0$ and $\mathcal S=\sum_jg_j$. Concavity, applied to the boosted reference and the solution, gives the barrier inequality
\begin{equation}\label{eq:barrier-alternative}
 -\sum_jg_j(u_t)_{j\bar j}
       \geq\xi+\varepsilon\mathcal S-Ng_k.
\end{equation}
The harmless additional $t\mathcal S\geq0$ may be discarded. If $Ng_k\leq\tfrac12\min\{\xi,\varepsilon\}(1+\mathcal S)$, the right side is at least $\tfrac12\min\{\xi,\varepsilon\}(1+\mathcal S)$. Otherwise the derivative in that direction is already comparably large. This is the form of the alternative used in Sun's estimate; it explains how a pointwise partial-phase inequality becomes a favorable term in a differential inequality. Sun carries it through the twisted-equation bounds. The old full-phase requirement is avoided because the limiting reference only sees the $n-1$ remaining angles, while the solutions themselves stay in the controlled full-phase cone.

\subsection{Openness, compatibility and the final limit}
The constant in \eqref{eq:twisted} is not an arbitrary extra datum. To see the compatibility in the linear problem, fix a solution $\beta_t=\alpha+t\omega+\ddc u_t$ and set
\[
 \Gamma_t=F_{n-1}(\beta_t,\omega;\theta),\qquad
 L_tv=\frac{n\,\ddc v\wedge\Gamma_t}{\omega^n}.
\]
The condition $P(\beta_t)<\theta$ makes $\Gamma_t$ a strictly positive $(n-1,n-1)$ form, as can be checked on each hyperplane by the sine formula used in Section~\ref{sec:numerical}. It is also closed. Thus $L_t$ is elliptic and Stokes' theorem gives
\begin{equation}\label{eq:linear-compatibility}
 \int_XL_tv\,\omega^n=0,\qquad
 \int_XvL_tv\,\omega^n
       =-n\int_X\ii\partial v\wedge\bar\partial v\wedge\Gamma_t.
\end{equation}
The kernel consists exactly of constants; self-adjoint elliptic theory gives the zero-mean range. On potentials normalized by $\int_Xv\omega^n=0$, adjoining a scalar $b$ makes $(v,b)\mapsto L_tv-b$ invertible. The implicit function theorem solves the augmented equation near $t$. Integrating that equation forces its scalar to be precisely the displayed $c_t$. In particular, one must not claim that an operator with constants in its kernel is invertible on all functions without treating normalization and range.

The estimates now make the path closed. Since $c_t\geq0$ and $D>0$, \eqref{eq:twisted-operator} implies $Q\leq\theta$ along it; bounded eigenvalues give a positive gap for every omitted arccotangent. A convergent sequence of normalized solutions therefore has a smooth limit in the required cone. At $t=0$, $c_0=0$ and $0<Q<\pi$ give $Q=\theta$, not $Q=\theta+\pi$. These two parts of continuity have different jobs: Fredholm compatibility gives a local family, while quantitative subsolution estimates prevent escape at its endpoint. Chen's twisted existence and Sun's estimate package remain explicit external inputs, not claims of a new complete proof here.

\section{From numerical positivity to analytic representatives}\label{sec:numerical}
\subsection{Necessity detects every complex tangent plane}
Suppose $Q_\omega(\alpha_u)=\theta\in(0,\pi)$. On a complex $p$-plane, let $\nu_1\leq\cdots\leq\nu_p$ be the eigenvalues of the restricted form and order the ambient eigenvalues increasingly. Eigenvalue interlacing gives $\nu_j\geq\lambda_j$. Since arccotangent is decreasing and positive,
\[
 0<q_V:=\sum_{j=1}^p\arccot\nu_j
 \leq\sum_{j=1}^p\arccot\lambda_j<\theta\quad(p<n).
\]
On this plane the density of $F_p$ relative to $\omega^p$ is
\begin{equation}\label{eq:sine}
 \prod_{j=1}^p\sqrt{1+\nu_j^2}
       \frac{\sin(\theta-q_V)}{\sin\theta}>0.
\end{equation}
Integration gives the proper-subvariety inequalities, including singular varieties by integration on their regular loci. Adding $t\omega$ decreases every arccotangent. For $t>0$ the full phase also lies strictly between $0$ and $\theta$, proving the top-dimensional ray inequality (in fact strictly for $t>0$). This elementary argument identifies what the numerical tests measure: positivity on all possible restricted phase configurations, not merely the sign of $Z$.

\subsection{Why projectivity eliminates the ray}
On a projective $X$, choose a rational ample class $\gamma$. Mixed intersections with powers of $\gamma$ are represented, up to positive constants, by effective hyperplane-section cycles. The assumed endpoint tests on their irreducible components imply positivity of
\[
 \int_V F_k(A,\Omega;\theta)\,\gamma^{p-k},\qquad 1\leq k<p.
\]
The term for $k=0$ is positive because $F_0=1$, and the term for $k=p$ is the endpoint test (or zero for $V=X$). Consequently the expansion
\begin{equation}\label{eq:ample-ray}
 \int_V F_p(A+t\gamma,\Omega;\theta)
   =\sum_{k=0}^p\binom pk t^{p-k}
       \int_V F_k(A,\Omega;\theta)\gamma^{p-k}
\end{equation}
is positive for $t>0$. The general test-family theorem of \cite{CLT} applies along this ample ray. It yields a solution, which in turn supplies the $\Omega$-ray tests by necessity. The argument uses a rational \emph{auxiliary} ample class, not rationality of the given K\"ahler class. On a nonprojective manifold these hyperplane cycles may not exist; endpoint positivity cannot then be substituted into the same proof.

\subsection{A wrong-branch example in dimension three}
Zhang \cite{Zhang} shows that the endpoint-only proposal fails on generic tori. The following explicit specialization makes both the branch and the missing test visible. Take a complex three-torus with no positive-dimensional proper analytic subvarieties, a flat K\"ahler form $\omega$, $\alpha=-\omega$, and
\[
 \theta=\pi/4,\qquad \hth=5\pi/4.
\]
All proper-subvariety tests are vacuous. Moreover
\[
 (-1+\ii)^3=2+2\ii,\qquad
 F_3(-\Omega,\Omega;\pi/4)=0,
\]
and $e^{-\ii\hth}\int_X(\omega-\ii\omega)^3>0$, so even the oriented integral phase is correct. Nevertheless at a maximum of any potential $u$, $\alpha_u\leq-\omega$, and therefore
\[
 Q_\omega(\alpha_u)\geq3\arccot(-1)=9\pi/4>\pi/4.
\]
No solution on the requested branch exists. The ray detects the defect at $t=2$, because
\[
 F_3(\Omega,\Omega;\pi/4)=-4\Omega^3<0.
\]
This is a failure of the general endpoint criterion, not an unresolved version of it, and not a counterexample to Theorem~\ref{thm:numerical}.

\subsection{Sufficiency: the object that has to be built}
Numerical inequalities do not initially produce a smooth form satisfying \eqref{eq:subsolution}. The difficult direction of \cite{CLT} constructs such an analytic representative through a strengthened induction on analytic subvarieties. Its interfaces explain why the theorem needs all subvarieties, not just divisors. The technical product-equation and regularization estimates are substantial imported results; the account below identifies their inputs and outputs.

For every proper analytic $Y\subset X$, the induction constructs a representative on a neighborhood of $Y$ whose \emph{full} phase satisfies $Q<\theta$ \cite[Theorem 3.1]{CLT}. Resolving singular $Y$ allows analytic estimates on a smooth space; small K\"ahler perturbations handle the degenerating pullback background. Local constructions near lower-dimensional singular or exceptional loci supply the induction data. This is a stronger neighborhood statement than merely having positive integrals on $Y$. It cannot hold with $Y=X$: a global $Q<\theta$ would give $\int_XF_n>0$, contrary to compatibility.

For the ambient step use the twisted path \eqref{eq:twisted}, whose constants are nonnegative by the ray hypothesis. Starting at large $t$, consider the potential first obstruction $t_0\geq0$ to continuing toward zero. To cross it one needs a subsolution in $A+t_0\Omega$, not a formal weak limit of solutions with uncontrolled eigenvalues. The construction passes to $X\times X$ and solves a twisted equation in a class of the form
\[
 \pi_1^*(A+t\Omega)+K\pi_2^*\Omega.
\]
Choose $\theta<\Theta_0<\pi$ and $K$ so that
\[
 \Theta_0-\theta<\frac{\pi-\Theta_0}{n},\qquad
 \cot\frac{\pi-\Theta_0}{n}<K<\cot(\Theta_0-\theta).
\]
These inequalities keep the enlarged product phases below $\pi$; they are not a decorative large-parameter assumption. The right-hand sides concentrate mass near the diagonal. Fiber integration of the resulting higher-degree forms returns closed $(1,1)$ forms in the desired class on $X$.

\subsection{What fiber integration preserves}
The untruncated fiber operation has a useful explicit check. Write $A_t=A+t\Omega$ and let $\Xi_s$ be the product solution in $\pi_1^*A_t+K\pi_2^*\Omega$. Suppressing the small resolution perturbations, which are unnecessary in the smooth ambient step, the operation in \cite[Section 5]{CLT} is
\begin{equation}\label{eq:fiber}
 \eta_s=\frac1{\mathcal V_K}
 \sum_{k=0}^{\lfloor(n-1)/2\rfloor}
 \frac{(-1)^k n!}{(n-2k)!(2k+1)!}
 (\pi_1)_*\bigl(\Xi_s^{n-2k}\wedge\pi_2^*\omega^{2k+1}\bigr),
\end{equation}
where $\mathcal V_K=\Imm(K+\ii)^n\int_X\omega^n>0$. Every summand has bidegree $(n+1,n+1)$ before integration over an $n$-dimensional fiber, so $\eta_s$ has bidegree $(1,1)$ and is closed.

Its class is not guessed from positivity. In the expansion of $[\Xi_s]^{n-2k}$ only the term with exactly one horizontal factor survives fiber integration. Its coefficient changes the factorial ratio in \eqref{eq:fiber} into $(-1)^k\binom n{2k+1}K^{n-2k-1}$. Hence
\[
 [\eta_s]=\frac{A_t\int_X\omega^n}{\mathcal V_K}
  \sum_k(-1)^k\binom n{2k+1}K^{n-2k-1}=A_t.
\]
This calculation explains why the specific polynomial, rather than a generic average of product forms, is used.

The phase estimate is a separate issue. Splitting $\Xi_s$ into horizontal, mixed and vertical blocks gives a corrected horizontal form $\widetilde\Xi_s$, analogous to a Schur complement, with
\[
 \eta_s=\mathcal V_K^{-1}(\pi_1)_*
 \bigl(\widetilde\Xi_s\wedge\Imm(\Xi_{s,V}+\ii\omega_V)^n\bigr).
\]
The product phase bound controls the sum of the corrected horizontal phase and the vertical phase. In particular, the vertical imaginary determinant supplies a \emph{positive} fiber measure of the fixed mass $\mathcal V_K$. Monotonicity and concavity of the appropriate cotangent-phase operator allow the estimate to survive this weighted averaging \cite[Lemma 5.2]{CLT}. The mixed-block correction is essential: the horizontal restriction alone is not the expression being averaged. Likewise, cohomological correctness alone would say nothing about preservation of a nonlinear cone.

\subsection{The positivity reserve and smooth gluing}

The mass concentration leaves a definite positive reserve. After truncation and passage to a limit, one obtains a current $T$ in
\[
 A+t_0\Omega-\delta\Omega\qquad(\delta>0)
\]
whose local regularizations satisfy, with a uniform margin,
\begin{equation}\label{eq:current-cone}
 P(T^{(r)})\leq\theta-\varepsilon_0,
 \qquad Q(T^{(r)})\leq\Theta_0-\varepsilon_0.
\end{equation}
Here the precise estimates use sufficiently small regularization scales and the local background comparisons of \cite[Theorem 5.4 and Section 7]{CLT}. This is not a claim that an arbitrary weak limit automatically satisfies a nonlinear current equation.

The reserve does not follow from ordinary weak convergence of \eqref{eq:fiber}. The mass-concentration estimate distinguishes a positive top-degree component carrying diagonal mass from lower-degree components without such concentration. Truncation isolates this positive contribution. The resulting estimate supplies $\delta>0$ in the smaller class above, even though the numerical hypothesis did not posit a single uniform lower bound over all subvarieties. One cannot obtain that bound by taking the minimum of an infinite collection of strict numerical inequalities; the product construction is what produces the needed analytic reserve.

A current is still not the promised subsolution. After adding a sufficiently large smooth background, its positive-Lelong-number level sets are proper analytic sets. The induction provides smooth local representatives near the relevant level sets. Away from them the local regularizations in \eqref{eq:current-cone} can be used. To see why the remaining gluing preserves admissibility, consider two smooth local potentials $a,b$ for the same reference $\alpha_0$. For a convex smooth approximation $h_\epsilon$ to absolute value, with $|h'_\epsilon|\leq1$, put
\[
 M_\epsilon(a,b)=\tfrac12\bigl(a+b+h_\epsilon(a-b)\bigr),\qquad
 w=\tfrac12(1+h'_\epsilon(a-b)).
\]
A direct differentiation gives
\begin{align*}
 \alpha_0+\ddc M_\epsilon(a,b)
   ={}&w(\alpha_0+\ddc a)+(1-w)(\alpha_0+\ddc b)\\
      &+\tfrac12h''_\epsilon(a-b)\,
                \ii\partial(a-b)\wedge\bar\partial(a-b).
\end{align*}
The last term is semipositive. Convexity of the admissible matrix sets controls the first two terms, and monotonicity under addition of semipositive forms controls the last one. Suitable separation of the potentials makes the regularized maximum agree with the desired local choice outside the overlap. In the actual construction the reserve $\delta\Omega$ and strict phase margins absorb the background and smoothing errors, so this local computation remains applicable; without a margin, it would not justify the argument.

The output is a \emph{global smooth} form in $A+t_0\Omega$ satisfying $P<\theta$ and $Q<\Theta_0$. The twisted existence theorem now crosses $t_0$, and continuation reaches zero. This is the precise stage where numerical positivity becomes the analytic input of Section~\ref{sec:analytic}.

Thus the proof has a closed chain: subvariety tests give neighborhood representatives; product mass concentration gives a current with a positivity reserve; controlled regularization and gluing give a smooth subsolution; elliptic continuity gives the solution. Replacing any of these outputs by the next one without explanation would hide the central content of numerical sufficiency.

\section{A variational construction: the tangent phase flow}\label{sec:flow}
\subsection{A metric and a gradient, with their domain}
In this section assume the stronger target range
\begin{equation}\label{eq:hyper}
 (n-1)\pi/2<\hth<n\pi/2.
\end{equation}
Let $\mathcal H$ consist of smooth potentials with $|\Th(\alpha_u)-\hth|<\pi/2$ everywhere. Set
\[
 \rho_u=\Ree\bigl(e^{-\ii\hth}(\omega+\ii\alpha_u)^n\bigr),\qquad
 \sigma_u=\Imm\bigl(e^{-\ii\hth}(\omega+\ii\alpha_u)^n\bigr).
\]
On this \emph{specified lift}, $\rho_u$ is a positive volume form. Its integral is $|Z|$, and $\int_X\sigma_u=0$. Define the complex energy
\[
 \mathcal E(u)=\frac1{n+1}\sum_{j=0}^n
       \int_X u(\omega+\ii\alpha_u)^j\wedge(\omega+\ii\alpha)^{n-j},
 \qquad \mathcal J=-\Imm(e^{-\ii\hth}\mathcal E).
\]
Closedness and integration by parts telescope its derivative:
\begin{equation}\label{eq:variation}
 d\mathcal E_u(v)=\int_Xv(\omega+\ii\alpha_u)^n,
 \qquad d\mathcal J_u(v)=-\int_Xv\sigma_u.
\end{equation}
With the positive metric $\langle v,w\rangle_u=\int_Xvw\rho_u$, the negative gradient of $\mathcal J$ is $\sigma_u/\rho_u$. Thus its flow is
\begin{equation}\label{eq:flow}
 \partial_tu=\tan(\Th(\alpha_u)-\hth)=:q,
 \qquad u(0)=u_0\in\mathcal H.
\end{equation}
It has the exact identities
\begin{equation}\label{eq:dissipation}
 \frac{d}{dt}\mathcal J(u)=-\int_Xq^2\rho_u,
 \qquad
 \frac{d}{dt}\Ree(e^{-\ii\hth}\mathcal E(u))=\int_X\sigma_u=0.
\end{equation}
The first is dissipation; the second controls the otherwise free additive constant when combined with oscillation estimates. Neither identity alone proves compactness of the flow.

\subsection{Global existence is not yet convergence}
Takahashi proves that under \eqref{eq:hyper}, every $u_0\in\mathcal H$ yields a global smooth solution of \eqref{eq:flow}; if a smooth $C$-subsolution exists, then $u(t)$ itself converges in $C^\infty$ as $t\to\infty$ \cite[Theorem 1.4]{Tak}. These are two different assertions. A subsolution is used for time-uniform estimates and convergence, not for the theorem's bare long-time existence. The phase of $u_0$ may be merely supercritical, since the lower edge of $\mathcal H$ is $\hth-\pi/2>(n-2)\pi/2$.

The key structural improvement is \emph{tan-concavity}: for fixed hypercritical target, the matrix function
\[
 F(\alpha_u)=\tan(\Th(\alpha_u)-\hth)
\]
is concave on this almost-calibrated domain \cite[Theorem 1.1]{Tak}. This supplies a concave parabolic operator even where the original phase is not concave. Its linearization in a diagonal frame has coefficients
\[
 F^{j\bar j}=\frac{\sec^2(\Th(\alpha_u)-\hth)}{1+\lambda_j^2}>0.
\]
The concavity has a visible algebraic mechanism. Write $z=\sum_j\arctan\lambda_j-\hth$ and, for a diagonal variation $v$, put $a_j=v_j/(1+\lambda_j^2)$. Direct differentiation yields
\[
 D^2F(v,v)=2\sec^2z\left[\tan z\left(\sum_ja_j\right)^2
                                      -\sum_j\lambda_ja_j^2\right].
\]
The phase inequalities under \eqref{eq:hyper} make the bracket nonpositive, including the case of one negative eigenvalue; this is the substantive algebra in \cite[Section 3]{Tak}. For off-diagonal variations the spectral second derivative contains the factors
\[
 \frac{F_i-F_j}{\lambda_i-\lambda_j}
   =-\sec^2z\frac{\lambda_i+\lambda_j}
                    {(1+\lambda_i^2)(1+\lambda_j^2)}\leq0,
\]
with the continuous interpretation at repeated eigenvalues. The sign follows because supercritical phase forces every pairwise eigenvalue sum to be positive. Thus the tangent change of speed supplies a compensating rank-one term in the diagonal Hessian and preserves the correct off-diagonal signs. It is an analytic improvement, not a mere reparametrization of the original flow in time.

Differentiating the flow gives a linear heat equation for the speed,
\begin{equation}\label{eq:speed}
 \partial_tq=F^{j\bar k}\partial_j\partial_{\bar k}q.
\end{equation}
The maximum principle bounds $q$ between its initial extrema. Therefore the phase stays a definite distance from the boundary of $\mathcal H$. This prevents loss of the chosen branch, but by itself does not bound large positive eigenvalues. The finite-time Hessian estimates in \cite[Section 4]{Tak}, followed by parabolic regularity, exclude a finite-time obstruction. Their constants may depend on the chosen finite time interval.

For convergence, the strict subsolution supplies bounded \emph{parabolic} level slices. Applied to the graph operator $F(\lambda)+\tau$, their normal alternative controls either the barrier term or all linearized directions. This step uses concavity of $F$, not just convexity of the elliptic phase levels. The parabolic contact-set argument first bounds the potential from below. The role of the conserved real energy in obtaining the upper bound can be made explicit. Change the reference within its class so that $u_0=0$, and put $\mathcal C(u)=\Ree(e^{-\ii\hth}\mathcal E(u))$. Then $\mathcal C(u(t))=0$. Hypercritical target makes $\mathcal H$ convex: its upper phase bound is automatic, and its lower bound is a convex phase superlevel condition. Thus the segment $su$, $0\leq s\leq1$, lies in $\mathcal H$, and the first variation gives
\[
 \nu_u:=\int_0^1\rho_{su}\,ds,\qquad
 \nu_u\geq c\omega^n,\qquad \int_X\nu_u=|Z|,
 \qquad \int_Xu\nu_u=\mathcal C(u)=0.
\]
The lower bound is uniform along the flow because the phase margin is preserved. As a consequence,
\[
 \int_Xu\omega^n=-c^{-1}\int_Xu(\nu_u-c\omega^n)
 \leq c^{-1}\bigl(|Z|-c\int_X\omega^n\bigr)(-\inf_Xu).
\]
The Green-function upper bound for $u$ turns this into $\sup_Xu\leq C(-\inf_Xu)+C'$. Together with the parabolic lower bound this gives the full, time-uniform $C^0$ estimate \cite[Proposition 2.4 and Lemma 5.3]{Tak}. Conservation therefore controls a normalization that energy dissipation alone would leave undetermined.

The largest-eigenvalue estimate and a spatial blow-up argument then give time-uniform gradient and Hessian bounds \cite[Section 5.2]{Tak}. The blow-up limit is again a bounded plurisubharmonic function on $\C^n$, so it cannot have the normalized nonzero gradient. Uniform ellipticity and parabolic regularity yield time-uniform higher estimates.

\subsection{Why the limiting speed is zero}
Even compactness plus decreasing energy would need an argument to identify a limit. Here \eqref{eq:speed} and the cohomological identity give a direct one. Uniform estimates make the heat equation uniformly parabolic with controlled coefficients. The parabolic Harnack argument in \cite[Section 5.3]{Tak} gives constants $C,c>0$ such that
\[
 \osc_Xq(t)\leq Ce^{-ct}.
\]
Since $\int_Xq\rho_u=\int_X\sigma_u=0$ with $\rho_u>0$, one has $\inf q\leq0\leq\sup q$, and hence $\|q(t)\|_{C^0}\leq Ce^{-ct}$. Integrating $\partial_tu=q$ proves uniform convergence of the actual potentials, without an arbitrary drifting normalization. The higher estimates upgrade it to smooth convergence. The limit satisfies $q=0$; inside $\mathcal H$ this is exactly \eqref{eq:equation}, not a different tangent branch. Uniqueness identifies its representative modulo the additive constant selected by the conserved real energy.

This construction explains the role of the variational structure without inserting a separate properness theorem. The original phase flow $\partial_tu=\Th(\alpha_u)-\hth$ and a cotangent-phase flow have different speeds and parabolic structures. Having the same stationary equation does not transfer the theorem's initial-data or convergence conclusions to them. No general convergence or current openness claim for those other flows is made here.

\section{Boundaries of the conclusions}\label{sec:boundary}
\subsection{Critical phase needs new uniform estimates}
A positive gap from $(n-2)\pi/2$ entered the elliptic estimates. It is therefore invalid to send that gap to zero in their constants without further work. A recent result does supply such work. In the convention of \eqref{eq:equation}, Fu--Yau--Zhang \cite[Theorems 1.1--1.2]{FYZ} prove the following critical statement: if
\[
 \int_X(\alpha+\ii\omega)^n\in\R_{<0}
\]
and a smooth $\underline u$ satisfies
\[
 \min_k\sum_{j\ne k}\arctan\lambda_j(\alpha_{\underline u})
                   >(n-3)\pi/2,
\]
then $\Th(\alpha_u)=(n-2)\pi/2$ has a unique smooth solution normalized by $\sup_Xu=0$. The nonzero negative-real integral and the \emph{strict} subsolution are part of the statement. This is the controlling version arXiv:2511.21492v4, dated 10 June 2026, treated here as a preprint rather than as a theorem whose entire proof has been independently certified.

The construction perturbs $\alpha$ to $\alpha+t\omega$. The compatible phases become supercritical with excess of order $t$, and the critical subsolution persists for small $t$. The new input is a $C^{2,\gamma}$ estimate independent of $t$, rather than the older estimate with constants depending on the inverse phase gap. It permits a smooth limit. Thus it is no longer accurate to describe the whole critical phase as lacking an existence theory, but it would be equally inaccurate to omit these hypotheses and claim unconditional critical solvability.

\subsection{Non-strict tests change the analytic problem}
At a numerical boundary, strict partial-phase gaps, positivity reserves for gluing, or smooth ellipticity can degenerate. The surface example already proves that non-strict intersection inequalities cannot generally imply a global smooth solution. To claim a weak replacement one must specify the potential class, the meaning of wedge products, the equation, the singular set and the convergence topology. Merely taking a weak limit of forms does not justify passing a nonlinear determinant to the limit.

No unrestricted weak-existence or weak-flow-convergence theorem is asserted here. Such results require their own hypotheses and solution concepts, outside the smooth chain developed in this article. Questions beyond that chain are not declared open merely because the selected theorems do not address them. The conclusions actually established by the selected sources are the smooth surface criterion, the strict supercritical analytic and numerical criteria, the conditional hypercritical tangent-flow construction, and the stated critical result. Their mechanisms show exactly where weakening a hypothesis requires a new argument.

\input{references.bbl}
\end{document}
